\problemname{Range Sum 2}
Du får givet en lista med $N$ heltal $A_0, A_1, \dots, A_{N-1}$.

Du måste hantera $Q$ stycken frågor av två typer:
\begin{itemize}
  \item Typ 1: givet heltalen $l$ och $r$, skriv ut summan av talen i intervallet $[l,r]$, det vill säga $A_l + A_{l+1} + \dots + A_r$.
  \item Typ 2: givet heltalen $x$ och $y$, sätt $A_x = y$.
\end{itemize}

\section*{Indata}
Den första raden innehåller två heltal $N, Q$ ($1 \leq N \leq 3 \cdot 10^5$, $1 \leq Q \leq 10^5$).

Nästa rad innehåller $N$ positiva heltal $A_0, A_1, \dots, A_{N-1}$ ($1 \leq A_i \leq 10^9$),
den givna listan.

Därefter följer $Q$ rader som vardera innehåller en fråga. Varje fråga börjar med talet $T$ ($T \in \{1,2\}$).

Om $T=1$ följer heltalen $l, r$ ($0 \leq l \leq r < N$).
I detta fall ska du skriva ut summan av talen i intervallet $[l,r]$.

Om $T=2$ följer heltalen $x$ och $y$ ($0 \leq x < N$, $1 \leq y \leq 10^9$).
I detta fall ska du sätta $A_x = y$.


\section*{Utdata}
För varje fråga med $T=1$, skriv ut summan i det efterfrågade intervallet.

\section*{Poängsättning}
Din lösning kommer att testas på en mängd testfallsgrupper.
För att få poäng för en grupp så måste du klara alla testfall i gruppen.

\noindent
\begin{tabular}{| l | l | p{12cm} |}
\hline
  \textbf{Grupp} & \textbf{Poäng} & \textbf{Gränser} \\ \hline
  $1$    & $50$       &  $N,Q \leq 1000$ \\ \hline
  $2$    & $50$       &  Inga ytterligare begränsningar. \\ \hline
\end{tabular}
